\documentclass[12pt]{article}
%\renewcommand\normalsize{\fontsize{10pt}{12pt}\selectfont}
\usepackage{graphicx} % Required for inserting images
\usepackage[left=2cm,right=2cm,top=2cm,bottom=1.5cm]{geometry}

\title{HW 5, part 1}
%\author{Dennis Wilkinson}
%\date{March 31 2025; due Monday April 7}

\begin{document}

\setlength{\parindent}{0pt}
\setlength{\parskip}{3mm}

FS 1 Wilkinson HW 5, due Monday April 28, 2025

1. We talked about the kinetic theory of a gas in class, where rebounding molecules are the source of pressure. This led to the Ideal Gas Law $PV=NkT$ as long as $\frac{3}{2}kT=\langle KE \rangle$, where $\langle KE \rangle$ is the average KE of a molecule in the gas. Since the Ideal Gas Law is experimentally confirmed, this is the correct relation between average KE and T in a gas.

a. Calculate the average KE of one air molecule in your room at a comfortable temperature (convert to Kelvin).

b. Air mostly consists of N$_2$ molecules. Find the ``average" (root mean square) speed $v_{rms}=\sqrt{\langle v^2 \rangle}$ of one of these molecules in m/s and mi/hr. This is a typical speed of an N$_2$ molecule in your room. You need to use kg for mass; one proton or neutron has mass $1.67 \times 10^{-27}$ kg.

c. Are the O$_2$ molecules in the air moving faster, the same speed, or slower than the N$_2$ molecules? Why?

d. Kinetic theory gives us a way to estimate the time between collisions of air molecules. Imagine a molecule which has a collision, then moves freely for time $\Delta t$ before having another collision. If the molecule's speed between collisions is $v$, then it moves through a length $v \Delta t$ between collisions. 

Say the molecule has a circular cross sectional area of radius $r$ (not that accurate); show that the volume it moves through between collisions is then $V=\pi r^2 v \Delta t$. Draw it out if it helps.

This volume of air must contain one molecule or so; otherwise collisions would be happening more frequently. So the number density $n$ of air must be around $1/V$ molecules/m$^3$.

Using $n=2.45 \times 10^{25}$ molecules/m$^3$, 0.3 nm for the the N$_2$ molecule radius, and the $v$ you calculated in (b), calculate $\Delta t$ between collisions. I got around 0.25 nanoseconds. 



2. Hot air balloons - the big ones people ride in - have a burner which the pilot uses from time to time to heat the air inside the balloon (it looks like a flamethrower pointing up into the balloon). Once the air is warmer, the balloon rises.

But it's weird because the balloon maintains basically the same volume $V$ all the time - the material is not stretchy. Also, it's open to the air at the bottom, so the pressure $P$ inside equalizes with the outside $P$ and doesn't change much. 

a. So what other quantity, besides $T$, changed in $PV=NkT$ when the pilot fires the burner? Does this quantity go up or down? How do you think it does this? (I noticed it when I went in a balloon.)

b. Explain why the process in part (a) makes the balloon float higher in the air. This is a buoyancy question, not really relevant to Monday's lecture.

3. On a hot day in Austin, the air conditioners are running hard. One typical AC unit has a pressure $P_1 = 105$ psi, which means pounds per square inch, in its \textit{vapor line}. This is the line which transports the working fluid (containing the heat it previously absorbed) from the house to the \textit{condenser}, which is the outside unit with the fan. In the condenser, the \textit{compressor} gives the fluid a much higher pressure $P_2 = 260$ psi. The fluid (now a liquid) now has a higher temperature $T_2$.

a. Say the temperature in the vapor line just before the compressor is $T_1= 70 ^{\circ}$F. Convert this into Kelvin.

b. When the compressor increases the pressure from $P_1$ to $P_2$, if the volume/density stayed the same (which it does not), find what the temperature in the liquid line would be using $PV=NkT$. Convert this to Fahrenheit. Hint: $N$ stays the same and we are assuming $V$ stays the same, so the ratio $P/T$ also stays the same.


You should get an answer which is far higher than the outside temperature - in fact this is too high for the AC to function correctly. It's because we erroneously assumed the density stayed the same. This is totally wrong - in fact, the gas becomes a liquid in the compressor.

c. The reason our number in (b) is off is that the density/volume of the coolant does not stay constant. Rather, $P$, $V$, and $T$ all change. Thus, we can't directly find $T_2$ with only $P_1$,  $P_2$ and $T_1$; we have 3 unknowns ($T_2$, $V_1$ and $V_2$) but only two $PV=NkT$ equations.

However, we can find $T_2$ if we assume the compression process is \textbf{adiabatic} - meaning, no heat is exchanged with the environment. This is a very accurate assumption here since the compression process is quick - there isn't time for heat to flow. 

In this case, $P_1V_1^\gamma=P_2V_2^\gamma$, where $\gamma = 1.12$ for a typical coolant fluid. Do a lot of algebra to show that in this case, 
\begin{equation}
    T_2 = T_1 \left( \frac{P_2}{P_1} \right)^{1-\frac{1}{\gamma}}. 
\end{equation}
Hint: start from $P_2V_2=NkT_2$, and eliminate $V_1$ and $V_2$ using the adiabatic equation $P_1V_1^\gamma=P_2V_2^\gamma$.

d. Calculate $T_2$ and convert back to $^{\circ}$F. It should be somewhat hotter than a typical Austin summer day.  

e. After the compressor, the hot fluid passes through the \textit{liquid line}, which is formed into coils inside the condenser (before going into the house) Google ``condenser coils" for images if you want. Heat flows out of the line into the outside air, aided by the fan. 

Is this process adiabatic? Why or why not? How would you compute the pressure in the liquid line after the condenser coils?


4*. Here we'll derive the adiabatic relation $P_1V_1^\gamma=P_2V_2^\gamma$. In class Thursday, we'll show that for an adiabatic process, $dU = PdV$ where $U$ is the internal energy and equals $\frac{1}{\gamma-1}NkT$. Gamma is known as the specific heat ratio (or heat capacity ratio) and we'll learn about it Thursday.

a. Show that $dU = PdV$ becomes 
\begin{equation}
    \frac{dP}{P} = - \gamma\frac{dV}{V}
\end{equation}
using $PV=NkT$, the chain rule on $dU$, and some algebra. 

b. Integrate both sides to get $PV^\gamma=$ constant. 

c. $\gamma$ is determined by the number of degrees of freedom in the molecules making up the gas. Read Feynman lecture 40, section 5 if you haven't learned about this before. Roughly many internal degrees of freedom does the coolant gas used in problem 3c have?

\end{document}